\section{Related work}
\label{sec:related}

This section covers the work this study builds on and the classes of work it deliberately declines.
Absence is a design decision here, not an omission, and it is recorded as one.

\subsection{Fish and Literature}
\label{sec:related-fish}

The published material on this game is human and informal, and neither of its two substantial
sources is quantitative or describes a program. McLeod's Pagat entry \cite{pagat} is the canonical
rules source and records several points on which dialects differ; ours is specified in
\S\ref{sec:game-dialect}. Develin's chapter \cite{develin} is the most detailed strategy writing we
located, describes the same dialect, and supplies two observations this design uses: an ask teaches
the opponents as much as it teaches a teammate, and a player may not declare a half-suit they hold
entirely. Secondary strategy pages \cite{dorsa,strategy-site,wikiliterature} restate blackballing,
asking the asker, and an endgame delegation heuristic; they are evidence about how the game is
played and talked about, not sources of quantitative claims.

\textbf{Computational prior art remains close to absent.} Searches across this project's cycles
found no academic work on this game, no public repository containing search, planning or a working
learner for it, and no coverage in RLCard \cite{rlcard}. The one agent located, Somani's
\texttt{literature} \cite{somani,somani-server}, plays the four-player, 48-card variant, which does
not raise the six-player allocation problem, and publishes no strength numbers; a closed-source
mobile application advertises bots with no stated methodology \cite{playstore}. Every novelty
statement in this paper therefore takes the form ``we found no prior published work that \ldots{}''
and should be read against that scope. One recent item is worth naming because it may become the
external calibration this line of work lacks: Valet packages twenty-one traditional
imperfect-information card games in a common description language with measured branching factors
and durations \cite{goadrich2026valet}, though whether any of the twenty-one is a team game is not
stated on its abstract page and we did not fetch the full text, so nothing here is compared against
it.

\subsection{Determinization, and why double-dummy analysis fails here}
\label{sec:related-determinization}

Determinization, or perfect-information Monte Carlo --- sample hidden cards consistently with the
public history, solve the resulting perfect-information game, and play the move with the best
average value --- is the standard recipe for hidden-hand card games \cite{levy,ginsberg-gib}, and is
unsound however many worlds are drawn. Frank and Basin named the two mechanisms: \emph{strategy
fusion}, in which the searcher plays a different strategy in each sampled world although a real
player must commit to one strategy per information set, and \emph{non-locality}, in which a node's
value depends on parts of the tree outside its own subtree
\cite{frank-basin-ai,frank-basin-aaai}. Long, Sturtevant, Buro and Furtak characterise when the
method nevertheless performs well, through leaf correlation, bias and a disambiguation factor
\cite{long-pimc}; Information Set Monte Carlo Tree Search replaces independent trees with a single
tree over information sets \cite{cowling-ismcts,goodman-ris}; the $\alpha\mu$ family backs up Pareto
fronts of world-vectors instead of scalar averages, repairing both mechanisms, and reports Bridge
3NT declarer play at 60.2\% to 63.0\% with twenty worlds \cite{cazenave2019alphamu}; and
knowledge-based paranoia search reasons about the worst case over deduced-consistent worlds rather
than the average, reporting over 1{,}000 points on the extended Seeger scale in Skat
\cite{edelkamp2021paranoia}.

Fish breaks the method for a reason none of that literature addresses, and
\S\ref{sec:domain-not} states it: the perfect-information game is degenerate rather than merely
fused, so $\alpha\mu$'s repair of unsoundness does not rescue the construction.

What this corpus adds is that the binding problem is not the abstract failure but \emph{the
statistic}. One determinization's value is the final half-suit differential, whose spread is an
order of magnitude larger than the differences between the leading candidates, so an argmax over
sampled means selects on residual noise --- and does so more confidently the more candidates it is
offered. That is the optimiser's curse rather than strategy fusion, and \S\ref{sec:arc-vsix-search}
reports it measured at \vsixCurseGap{}~points by a control that forces the same code to return the
blueprint's own pick. The transferable claim from that result is the guard, not the search.

\subsection{Search in imperfect-information and cooperative games}
\label{sec:related-search}

States in an imperfect-information game do not have well-defined values, so depth-limited search
needs something at the leaf that represents the opponent's ability to choose. Brown, Sandholm and
Amos supply it: let the opponent select among $k$ continuation strategies at the depth limit, each
yielding its own leaf-value vector, and require the searcher to be good against that choice
\cite{brown2018depthlimited}. The soundness argument is stated for two-player zero-sum games, and
Fish's team-level game presents three uncorrelated opponent seats at the leaf, so nothing is
inherited; it is adopted here as a robustification heuristic with a measurable consequence, and
audited rather than assumed.

Two lines address what a search may prune. Zhang and Sandholm solve subgames over low-order
knowledge rather than the full common-knowledge closure, and prove that the naive form can
\emph{increase} exploitability \cite{zhang2021subgamesolvingwithoutck}. That warning applies directly
to the live-ask gate this policy family ships, which is a knowledge-limited pruning rule, and it is
why the side-channel and pathology gates of \S\ref{sec:methodology-gate} were built before any
search work rather than after it. Subgame solving in adversarial team games builds a gadget over the
team belief DAG and runs in polynomial time given a best-response oracle \emph{when the blueprint is
sparse} \cite{zhang2022teamsubgame}; a Fish blueprint is not sparse, so that escape hatch does not
open here.

In cooperative games the constraint transfers verbatim. SPARTA improves a blueprint purely at test
time and distinguishes single-agent search --- which assumes every other agent strictly follows the
blueprint --- from multi-agent search, in which all agents run the \emph{same} common-knowledge
procedure \cite{lerer2020sparta}. The moment a second teammate also searches, the first teammate's
belief over its partner is computed under a policy that partner is no longer playing, and the
improvement guarantee is void. Fish satisfies the requirement cheaply, for the reason
\S\ref{sec:domain-team} gives; the qualifier about randomisation is what
\S\ref{sec:arc-vsix-ties} turns into an engine switch.

The full-solver line \cite{moravcik2017deepstack,brown2020rebel,schmid2023sog} established that
test-time re-solving over public belief states is principled rather than a hack. The algorithms are
declined here, because each requires iterating or learning over a closure whose cardinality at deal
time is $45!/(9!)^5 \approx 1.9 \times 10^{28}$ deals from one seat. Update equivalence recasts
search as replicating a last-iterate-convergent update locally, needing no public-information
enumeration \cite{sokota2024updateequivalence}; its improvement guarantee is for common-payoff
games, which Fish is not. Learned Belief Search \cite{hu2021lbs} is redundant here, because this
game's exact belief is already computed and oracle-validated. Online planning by policy fine-tuning
\cite{fickinger2021rlfinetuning} requires gradient updates at decision time and does not transfer to
a CPU-only engine with no network. The depth-limited counter-strategy line
\cite{milec2021cdr,milec2025abd} carries one warning this study honours --- robust-response methods
do not compose with limited-look-ahead solving off the shelf --- and the learned-model and
gadget-refinement work \cite{kubicek2025lamir,kubicek2026refinements} is recorded for its
meta-lesson rather than its machinery: when a search subproblem has many near-equivalent optima,
which one is selected is not a tie-break. That is the same observation \S\ref{sec:arc-vsix-ties}
makes about an argmax resolved by array order.

Every method above is published with a caveat that it can increase exploitability, so none can be
evaluated by win rate alone. Local best response gives a cheap lower bound: a small greedy search,
at most two actions deep, showed every entrant in a 2016 poker competition to be more than
3{,}180\,mBB/h exploitable \cite{lisy2017lbr}. Its two documented pitfalls are carried into
\S\ref{sec:exploit} --- restricting where the response may begin deviating understates
exploitability, and a negative score proves nothing --- and every exploitability figure in this
paper is a lower bound within its response class.

\subsection{Policy-aware inference}
\label{sec:related-inference}

Skat engines enumerate the deals consistent with the public record and reweight them by a model of
what each player would have done --- a supervised per-card location model
\cite{solinas-supervised}, or reach probabilities taken from a policy \cite{rebstock-inference} ---
and a dedicated inference module was reported worth $+217$ tournament points per 36 games to the
engine Kermit on its own \cite{buro-kermit}. The policy prior of \S\ref{sec:engine-prior} is the
same idea in a much simpler form: a fitted soft reweighting of the rules posterior, with two
parameters and no learned opponent model.

This is the only literature that has measured the same trade-off, and it is where
\S\ref{sec:arc-vsix-belief}'s central negative result belongs. What is new there is the margin and
its direction: in this game the exact conditional posterior is computable in closed form and cheap,
so the usual excuse for approximating does not apply --- and yet substituting the exact posterior
into an otherwise identical policy \emph{loses}, on two independent seed banks.

Two adjacent lines are named and not used. Belief modelling in cooperative games --- the Bayesian
Action Decoder's public-belief formulation \cite{foerster2019bad} and its factorised-prescription
successor \cite{sokota2021capi} --- is aimed at games whose exact belief is intractable, which is the
opposite of the situation here. And the exact-counting machinery this study does not need to
approximate --- Ryser's inclusion--exclusion formula \cite{ryser}, Glynn's sign-vector formula
\cite{glynn}, the Jerrum--Sinclair--Vigoda chain \cite{jsv}, matrix scaling \cite{lsw} and
sequential importance sampling \cite{chen-sis} --- is sidestepped entirely by the block structure of
\S\ref{sec:engine-belief}. No learned belief model appears anywhere in this corpus, and that is a
stated non-goal rather than an oversight.

\subsection{Team games and common information}
\label{sec:related-team}

The formal statement of three teammates sharing a payoff, holding different information and unable
to communicate is the common-information approach \cite{nayyar2013common}, discussed in
\S\ref{sec:domain-team}. The team-game equilibrium literature is cited to delimit scope: the target
concept for a Fish team is a team-maxmin equilibrium with correlation \cite{celli-gatti}, connected
to extensive-form correlation by Farina and coauthors \cite{farina-collusion} and made computable
through the team belief DAG and the two-player conversion \cite{zhang-tbdag,carminati-tpi}, with
subsequent work on hidden-role settings \cite{carminati2024hiddenrole}. None of it is computed here.
The deeper reason a Fish team resists the DAG is that its meta-player fails A-loss recall, for
exactly one reason: each seat privately observes nine cards its teammates do not.

The cooperative-play line is cited for its evaluation discipline rather than its algorithms. Hanabi
was proposed as a benchmark precisely because self-play scores overstate coordination
\cite{bard-hanabi}, and the ad-hoc-teamplay results that followed
\cite{hu-sad,hu-otherplay,hu2021obl,hu2022sed,cox-hanabi,tian2020jps,jacob2022pikl,hu2022humanregularized,dizdarevic2025ah2ac2}
are the reason the cross-play measurement of \S\ref{sec:results-crossplay} exists at all: a
convention that collapses when partners are swapped is brittleness, and it is invisible to self-play
evaluation. The same lesson from a different direction is Cicero's use of a human-regularised policy
\cite{fair2022cicero}.

\subsection{Evaluation methodology}
\label{sec:related-eval}

The apparatus of \S\ref{sec:methodology} draws on a literature that is usually treated as
plumbing. Duplicate-style variance reduction and its estimators --- DIVAT \cite{zinkevich-divat} and
AIVAT \cite{burch-aivat} --- are the reason every cell in this paper is a paired duplicate rather
than an independent sample. The Annual Computer Poker Competition's protocols \cite{acpc} are the
model for a fixed opponent panel scored identically for every arm. Abstraction pathology
\cite{waugh-abstraction} is the standing warning that a model can be improved on its own metric and
made worse in play --- which \S\ref{sec:rejected-leaf} reproduces exactly. Re-evaluation of agent
populations \cite{balduzzi-reeval} and the rating literature
\cite{coulom-whr,herbrich-trueskill,nelo} inform why this paper reports a worst case and a minimax
regret over a shared panel rather than a single rating number. Recent work on gold-standard
evaluation practice \cite{kelidari2026goldstandard} is consistent with the preregistration and
sealing discipline of \S\ref{sec:methodology-seal}, which this project adopted independently.
